\section*{Lời mở đầu}
\addcontentsline{toc}{section}{Lời mở đầu}

Tại nhiều nút giao đô thị, đèn tín hiệu vẫn hoạt động theo chu kỳ cố định. Cách điều khiển này đơn giản và dễ triển khai, nhưng chưa phản ánh đúng sự thay đổi liên tục của mật độ phương tiện theo từng hướng. Khi một hướng đông xe nhưng thời gian xanh ngắn, trong khi hướng còn lại ít xe vẫn được cấp đèn xanh dài, hệ thống tạo ra thời gian chờ không cần thiết và làm giảm hiệu quả khai thác mặt đường.

Báo cáo này trình bày quá trình xây dựng mô hình thử nghiệm \projectname{}: hệ thống đèn giao thông thích nghi kết hợp thị giác máy tính và bộ điều khiển nhúng. Trong đó, camera hoặc video cung cấp dữ liệu giao thông; máy tính biên xử lý phát hiện phương tiện, chia hướng theo vùng quan tâm, ước lượng mật độ và tạo kế hoạch thời gian đèn; vi điều khiển nhận kế hoạch qua UART rồi tự vận hành máy trạng thái điều khiển đèn đỏ, vàng, xanh.

Vì đây là báo cáo môn Vi xử lý, trọng tâm được đặt trước hết vào phần điều khiển nhúng: phân chia trách nhiệm giữa host và MCU, giao thức UART, hàng đợi FreeRTOS, task, mutex, task notification, watchdog và máy trạng thái đèn giao thông. Phần trí tuệ nhân tạo được trình bày sau với vai trò tạo dữ liệu điều khiển mức cao. Cách phân bổ này phản ánh đúng bản chất của dự án: đây là một hệ thống nhúng có AI hỗ trợ, không phải chỉ là một bài toán nhận diện đối tượng.

Báo cáo không khẳng định các mức giảm ùn tắc nếu chưa có thí nghiệm đối chứng đầy đủ. Các kết quả được trình bày chủ yếu là kết quả triển khai phần mềm, kiểm thử pipeline, kiểm thử runtime, benchmark FPS và bằng chứng cho việc bộ điều khiển không đổi đèn tùy tiện theo từng khung hình.

